% NAL 2338, batch 6 — Gallica views 101–104, the last four of the volume.
% Pass: Opus 5.5 (claude-opus-5-5), 2026-09-23 — first pass, unchecked against the views by a human.
%
% What the four views hold. Greyscale scans (v. 104 in colour, read in grey),
% portrait, about 4234–4316 × 5740, one page per view, the leaf lying on the
% opened volume with the edge of the next leaf or of the guards at one side.
% Three views carry text, all in one hand and all one piece: the end of a
% Latin letter on experiments with quicksilver, water and wine in glass tubes
% and siphons, which names « D. de Paschal » as the one who made them.
%
%   v. 101     Recto numbered « 49 ». Continues the text of view 100 without
%              a break (see Continuity). The writer's opponents hold that the
%              space at the top of the tube is filled with « spiritus »; one
%              of them, « Scriptor ille », has published within twenty-four
%              hours a « libellum » « De Mirabilibus Mercurij (sic enim
%              hydrargyrum vocant) proprietatibus » (so the leaf; no author
%              or printed book is named or identified here). « D. de
%              Paschal » has glass tubes of 40 feet tied to a mast, invites
%              everyone to a large place « in area officinæ Vitriariorum »,
%              and, having computed that the 2 7/24 feet of quicksilver are
%              balanced by about 31 1/9 feet of water and 31 2/3 of wine, gets
%              his opponents to concede that wine, having more « spiritus »,
%              should leave more space at the top; the mast is raised, the
%              water stands at 31 1/9 and the wine at 31 2/3; the liquids are
%              exchanged, the mast inclined; a tube of 15 feet partly of water,
%              partly of quicksilver. The sentence runs on.
%   v. 102     Verso, unnumbered. The same experiment made in tubes of 4 or 5
%              feet; the ratio of the weight of water to that of quicksilver
%              « vt 1. ad 13. 4/7 »; the experiments « quater publice
%              repetitis »; what the apparent vacuum is: animals to be put in
%              tubes to see whether they live and fly; the « materia
%              subtilissima » of others refused; « Quidquid sit », at least
%              the vacuum of the Peripatetic elements. Then « D. de Paschal »'s
%              experiment with bent tubes, « Syphones », their two mouths in
%              two dishes of quicksilver. The sentence runs on.
%   v. 103     Recto numbered « 50 ». End of the siphon experiment (2 7/24 feet
%              above each dish; the « venter » of the siphon appears empty);
%              no siphon or suction pump will raise water above 31 feet; « His
%              Animaduersis », two questions debated « inter eruditos »
%              (Prima: how the liquid is held balanced at its height; Secunda:
%              why a drop of another liquid let in makes the column rush to
%              the top, breaking tubes); the writer will wait for his
%              correspondent's letters « super hac re ». Closing: « Vale Vir
%              mihi Colendiss.me et me ut hucusque fecisti in posterum amare
%              perge. », then, alone at the foot, the date « Parisijs 12
%              Calend. Octob. 1647 », the year underlined.
%   v. 104     Verso of the leaf numbered 50: blank. The text of v. 103 shows
%              through reversed (the date legible in mirror at the foot);
%              checked at contrast and mirrored, it matches v. 103 line for
%              line and carries nothing of its own. Not transcribed. Last view
%              of the volume.
%
% Continuity. View 101 CONTINUES a text begun before it: the last line of
% view 100 (read on the mirror only to settle this) ends « Sicque Hydrargirum
% in tubis quidem aliis, altius; », and view 101 opens « in aliis autem
% humilius stare debuisset ». View 100 is in the same hand and layout; it
% belongs to batch 5, and where the letter begins is for the coordinator to
% reconcile. Views 101 → 102 → 103 run on in mid-sentence. The letter ends on
% view 103 with its valediction and date; view 104 is blank, and nothing runs
% on (it is the last view of the volume).
%
% Who writes to whom. A letter in the first person (« ut jam dixi »,
% « expectabo … Epistolas tuas », « Vale Vir mihi Colendiss.me »), to a
% correspondent addressed in the second person. On these four views the
% letter carries no salutation, no signature and no address: neither writer
% nor addressee is named. « D. de Paschal » (D.nus, D.no) is named in the third
% person as the maker of the experiments. Nothing on these leaves attributes
% the hand or the letter to Mersenne. The only date is the leaf's own,
% « Parisijs 12 Calend. Octob. 1647 », transcribed as written and not
% converted.
%
% One hand. A posed, regular sloping italic, a fair copy almost without
% correction (a handful of deletions and interlinear substitutions: « de »
% → « cum », v. 101; « quibus » → « quos », and a line recast, v. 102;
% « descendat » → « decedat », and « firmiss.mi » added in the margin,
% v. 103), with long line-fillers at the line ends (not transcribed). It
% writes v at the head of a word and u inside it (« vnus », « vtraque »,
% « seruare », « grauitatem », « obseruatio »), kept. The ligature æ is
% written as a with a small stroke above (« aquā »), set æ; « equalitatem »
% (v. 101) and « aquali » (for æquali, v. 102) are written without it. -que
% as q; kept. -ij kept (« Mercurij », « Parisijs »). The numeral 5 is a long
% s, so that « 15 » looks like « is »; the 7 has a hooked form that can be
% taken for a 4 (the 4 is closed). A long s is transcribed as s.
%
% Notation. Mixed numbers written with a stacked fraction after the integer
% (2 7/24, 31 1/9, 31 2/3, 13 4/7), set $2\frac{7}{24}$ etc. The writer's
% « 13. 4/7 » with a point is kept. No figure.
%
% Numbers on the leaves.
% - « 49 » (v. 101) and « 50 » (v. 103), in ink at the top right of the
%   rectos; the versos (v. 102, 104) carry none. They continue the thin ink
%   series that batch 1 found on every recto across all hands (1 at v. 2 …),
%   and « 50 » on the last written leaf matches the library's certificate
%   (« Volume de 50 Feuillets / 1er Octobre 1888 », v. 1). It behaves as the
%   library's foliation, but it is ink, not pencil, so, as for the rest of
%   this volume and for NAF 5176, NAF 9544 and Français 12357, no \folio{} is
%   written until a human has looked.
% - A few ink dots in the left margin of v. 101 (at « autem pedes 31 2/3 »):
%   no figure.
% - The round stamp « BIBLIOTHEQUE NATIONALE … R.F. » on v. 103.
% - All other figures are content (feet, the ratio 1 : 13 4/7).
%
% The catalogue gives no dating: \dating{} is left empty, and this pass
% infers no date for the volume. The date above is the leaf's.

% Coordinator's reconciliation (2026-09-23), after all six batches:
%   One ink series of leaf numbers runs top right of the rectos across every
%   hand and piece: 1–10 (v. 2–20), 11–20 (v. 22–40), 21–29 (v. 42–59),
%   30–38 with « 34 bis » for the figure slip (v. 61–79), 39–48 (v. 81–99),
%   49–50 (v. 101–103). It ends on « 50 » at the last written leaf, as the
%   certificate of v. 1 has it (« Volume de 50 Feuillets / 1er Octobre
%   1888 »): it is the library's count. Being ink, it gets no \folio{}, in any
%   batch. The bold series on the letter of v. 24–35 (1–11) is that letter's
%   own pagination (batch 2).
%   Seams: v. 60 ends « vniuersa- » and v. 61 opens « lius »; batch 3 had set
%   the catchword into v. 60, and now ends where the page does. v. 100 runs
%   into v. 101 (« altius ; » / « in aliis autem humilius »), the « De Vacuo
%   Narratio » (v. 97–103). The « non » sign is set « nl » in batch 1 (hand B,
%   an n with a barred ascender) and « n^o » in batch 2 (the fast cursive, an
%   n with a raised loop): two hands, two signs, left as each pass read them.

\documentclass[11pt,a4paper]{article}
% The preamble shared by every transcription of the mathematicians' manuscripts in Gallica.
%
% It rests on one idea: **uncertainty compounds, it does not resolve.** A
% transcription that smooths over a doubtful word has destroyed the only thing
% it was made for — a reader must be able to tell, at every point, what was on
% the page and what was supplied. The apparatus macros below are therefore the
% critical apparatus, not decoration.
%
% The same commands are recognised by `scripts/render.mjs`, which produces the
% site's reading view, and by `scripts/tei.mjs`. A macro added here must be
% added there too, or rendering fails loudly — which is the wanted behaviour.
%
% Comments here are English, like the rest of the repository. The documents
% this preamble sets are French, and that is deliberate: see the skills.

\ProvidesPackage{germain}[2026/09/15 Transcriptions of mathematicians' manuscripts in Gallica]

\usepackage{amsmath,amssymb,amsthm}
\usepackage{mathrsfs}
% Kept from the parent preamble so the renderer's diagram support stays
% available; these manuscripts draw no commutative diagrams, and nothing here uses it.
\usepackage{tikz-cd}
\usepackage[english,french]{babel}
\usepackage{fontspec}

% --- Greek ------------------------------------------------------------------
% The text face stays fontspec's default, Latin Modern, which has no Greek:
% XeTeX drops every glyph it lacks with a line in the log and nothing on the
% page, and the Greek words the manuscripts quote vanished from the PDFs. So
% the Greek and Greek Extended blocks get a XeTeX character class of their own,
% and entering or leaving a run of it switches the family to CMU Serif — the
% same Computer Modern design, with polytonic Greek — and back. Only the
% family changes: a Greek word in \textit or \textbf keeps its shape and series.
% The transitions to and from every other class carry the tokens that class
% has towards an ordinary letter, so French spacing before « ; » or inside
% « » still holds after a Greek word. No grouping, so no brace or \endgroup
% can come out unbalanced; maths is left alone, since XeTeX inserts nothing
% there. Kept identical in `exercices.sty`.
\makeatletter
\newfontfamily\germainGreekfont{cmun}[
  Extension=.otf, NFSSFamily=germaingreek,
  UprightFont=*rm, ItalicFont=*ti, SlantedFont=*sl,
  BoldFont=*bx, BoldItalicFont=*bi, BoldSlantedFont=*bl]
\newXeTeXintercharclass\germain@greek
\def\germain@greekrange#1#2{%
  \@tempcnta=#1\relax
  \loop
    \XeTeXcharclass\@tempcnta=\germain@greek
    \ifnum\@tempcnta<#2\advance\@tempcnta\@ne\repeat}
\germain@greekrange{"0370}{"03FF}
\germain@greekrange{"1F00}{"1FFF}
\def\germain@greekprev{\rmdefault}
\def\germain@greekin{%
  \edef\germain@greekprev{\f@family}\fontfamily{germaingreek}\selectfont}
\def\germain@greekout{\fontfamily{\germain@greekprev}\selectfont}
\def\germain@greekjoin{%
  \XeTeXinterchartokenstate=\@ne
  \@tempcnta=\z@
  \loop
    \ifnum\@tempcnta=\germain@greek\else
      \XeTeXinterchartoks\germain@greek\@tempcnta=\expandafter{%
        \expandafter\germain@greekout\the\XeTeXinterchartoks\z@\@tempcnta}%
      \XeTeXinterchartoks\@tempcnta\germain@greek=\expandafter{%
        \the\XeTeXinterchartoks\@tempcnta\z@\germain@greekin}%
    \fi
    \ifnum\@tempcnta<4095\advance\@tempcnta\@ne\repeat}
% babel-french sets its own classes, and saves and restores the interchar
% state, each time a language is selected: join again after it does.
\addto\extrasfrench{\germain@greekjoin}
\addto\extrasenglish{\germain@greekjoin}
\AtBeginDocument{\germain@greekjoin}
\makeatother

\usepackage{geometry}
\geometry{a4paper,margin=25mm,marginparwidth=28mm}
\usepackage{marginnote}
\usepackage{xcolor}
\usepackage[normalem]{ulem}
\setlength{\emergencystretch}{3em}
\usepackage{hyperref}
\hypersetup{colorlinks=true,linkcolor=black,urlcolor=black,pdfborder={0 0 0}}

% --- Batch metadata ---------------------------------------------------------
% Taken as they stand by the rendering script, which shows them at the head of
% the reading view. They fix what the file claims to cover. The macro names are
% the parent project's, so that the scripts stay shared; \folder here means the
% volume's slug — `fr-9115` — and \pages counts Gallica views.

\newcommand{\folder}[1]{\gdef\grothFolder{#1}}
\newcommand{\batch}[1]{\gdef\grothBatch{#1}}
\newcommand{\pages}[2]{\gdef\grothFirst{#1}\gdef\grothLast{#2}}
\newcommand{\foldertitle}[1]{\gdef\grothFoldertitle{#1}}

% The holder's own shelfmark, printed as they write it: « Français 9115 ».
\newcommand{\shelfmark}[1]{\gdef\grothShelfmark{#1}}

% The dating, copied verbatim from the holder's catalogue — never inferred
% here. « XIXe siècle » is the BnF's; a pass that dates a leaf from its content
% says so in a \note{}, as its own inference.
\newcommand{\dating}[1]{\gdef\grothDating{#1}}

% The status notice, printed in the title block and repeated as a diagonal
% watermark on every page of the PDF. The manuscripts are in the public
% domain; what this line declares is not a right but a quality — an
% unverified first pass by a machine. The skills write the line; a file
% without it gets no watermark, so leaving it out is visible in the output.
\newcommand{\watermark}[1]{\gdef\grothWatermark{#1}}

\gdef\grothFolder{?}\gdef\grothBatch{}\gdef\grothFirst{?}\gdef\grothLast{?}%
\gdef\grothFoldertitle{}\gdef\grothShelfmark{}\gdef\grothDating{}\gdef\grothWatermark{}

% --- The résumé -------------------------------------------------------------
\newenvironment{resume}
  {\small\setlength{\parskip}{0.45em}}
  {\par\medskip\hrule\medskip}

% --- Keywords ---------------------------------------------------------------
% In English, closing the modernised reading's résumé: the modern vocabulary
% under which to search for what the leaves construct. The manifest extracts
% them to tag the volume — this line is the single source of the tags.
\newcommand{\keywords}[1]{\par\smallskip\noindent
  {\footnotesize\textsc{Keywords} --- \textit{#1}}\par}

% --- The critical apparatus -------------------------------------------------

% The Gallica view number, in the margin. This is the anchor that lets a
% disagreement over a formula be settled by looking at *one* image rather than
% a volume of 750. The site uses it to turn the facsimile as the transcription
% is scrolled. A view is one scanned image: usually one side of a leaf.
\newcommand{\page}[1]{%
  \marginnote{\footnotesize\color{gray!70}\textsf{v.\,#1}}%
  \label{page:#1}\ignorespaces}

% The BnF's pencilled foliation, where the leaf carries it and a pass can read
% it: « 198r ». This is what the literature cites — « ff. 198r–208v » — and
% the one line that turns a citation into a view. Recorded, never inferred:
% a folio a pass cannot read is not written.
\newcommand{\folio}[1]{%
  \marginnote{\footnotesize\color{gray!50}\textsf{f.\,#1}}\ignorespaces}

% The modernised reading's coarser counterpart to \page{}: which views a
% section is drawn from.
\newcommand{\pagerange}[2]{%
  \marginnote{\footnotesize\color{gray!70}\textsf{v.\,#1--#2}}%
  \ignorespaces}

% Illegible: never guessed.
\newcommand{\ill}{\textcolor{red!60!black}{[\,\ldots\,]}}

% A doubtful reading: the word is given, and flagged as uncertain.
\newcommand{\uncertain}[1]{\uline{#1}}

% An editorial addition — a missing letter, an implied word.
\newcommand{\add}[1]{\textcolor{blue!50!black}{[#1]}}

% Struck out by the writer. What was crossed out is part of the document.
\newcommand{\struck}[1]{\sout{#1}}

% The transcriber's note, on the state of the page or the mathematics.
\newcommand{\note}[1]{\par\smallskip{\small\color{gray!60!black}\textit{[#1]}}\par\smallskip}

% A marginal note of hers, distinct from the body of the text.
\newcommand{\marginal}[1]{\par\smallskip\noindent\hspace{1em}%
  {\small\color{gray!60!black}#1}\par\smallskip}

% --- Title ------------------------------------------------------------------
% The title block is French, like the documents themselves.

\AddToHook{shipout/background}{%
  \ifx\grothWatermark\empty\else
    \begin{tikzpicture}[remember picture, overlay]
      \node[rotate=52, scale=3.4, align=center, text=gray!18,
            font=\sffamily\bfseries]
        at (current page.center) {\grothWatermark};
    \end{tikzpicture}%
  \fi}

\AtBeginDocument{%
  \begin{center}
    {\large\bfseries Manuscrits de mathématiciens (Gallica) ---
      \ifx\grothShelfmark\empty\grothFolder\else\grothShelfmark\fi}\\[2pt]
    {\itshape\grothFoldertitle}\\[4pt]
    {\small vues \grothFirst--\grothLast
      \ifx\grothBatch\empty\else\ (lot \grothBatch)\fi}\\[2pt]
    \ifx\grothDating\empty\else
      {\small\color{gray!60!black}Datation du catalogue : \grothDating}\\[2pt]
    \fi
    \ifx\grothWatermark\empty\else
      {\small\bfseries\color{red!45!gray}\grothWatermark}\\[2pt]
    \fi
    {\footnotesize\color{gray!60!black}Transcription établie d'après les images IIIF de Gallica.
     Source gallica.bnf.fr / Bibliothèque nationale de France.\\
     Les lectures incertaines sont soulignées, les illisibles notées [\,\ldots\,],
     les ajouts entre crochets ; \textsf{v.} numérote les vues de Gallica,
     \textsf{f.} la foliotation au crayon de la BnF.}
  \end{center}
  \vspace{1em}}

\endinput


\folder{nal-2338}
\shelfmark{NAL 2338}
\batch{6}
\pages{101}{104}
\dating{}
\watermark{Lecture automatique\\non vérifiée}
\foldertitle{Papiers divers de Mersenne.}

\begin{document}

\page{101}
\note{Recto. En haut à droite, à l'encre, « 49 » (série continue ; voir l'en-tête). La page continue sans rupture le texte de la vue 100, dont la dernière ligne s'arrête sur « Sicque Hydrargirum in tubis quidem aliis, altius; » : « in aliis autem humilius » en est la suite. Même main, italique posée et régulière ; longs traits de remplissage en fin de ligne, non transcrits.}

in aliis autem humilius stare debuisset, nec seruare equalitatem illam pedum $2\frac{7}{24}$. Cum tamen illi obstinatius opinionem suam extollerent, ut jam eorum vnus quem \uncertain{exteri} veluti ducem sequebantur, libellum 24 horarum spatio editum, in lucem emisisset, De Mirabilibus Mercurij (sic enim hydrargyrum vocant) proprietatibus ; censuit D.\textsuperscript{nus} De Paschal alia ratione eis occurrendum esse. Curauit igitur duci ex Crystallo Tubos 40 pedum, eosque malo alligari, et machinas instrui, ut jam dixi ; constitutoque die, ac loco amplissimo in area officinæ Vitriariorum, inuitauit omnes ut adessent, mira conspecturi. Adfuerunt atque una cum illis Scriptor ille suis asseclis vndique stipatus.
\note{le numérateur de $2\frac{7}{24}$ est tracé ici d'un crochet qui le rapproche d'un 4 ; il est net à la ligne 31 de la page (« pedum 2 7/24 solius hydrargyri »), et c'est la même hauteur (« equalitatem illam »). « exteri » : c'est ce que porte la page, ex- bien formé ; le sens attend plutôt « cæteri ». Le début de « proprietatibus » est traversé par le trait qui prolonge la parenthèse fermante ; le mot n'est pas lu comme biffé.}

Inierat autem priuatim solertissimus D. de Paschal, calculum de aqua et vino \struck{de} \add{cum} hydrargyro comparatis secundum grauitatem, ut inde debitam vnicuiq; altitudinem eliceret, ad hoc ut in iis altitudinibus æquiponderarent : repereratque posita altitudine prædicta hydrargyri pedum $2\frac{7}{24}$ deberi aquæ pedes $31\frac{1}{9}$ circiter : Vino autem pedes $31\frac{2}{3}$ proxime. Itaque antequam quidquam sui propositi aperiret, interrogauit egregios illos Sapientes nec difficulter ab iis elicuit, Majorem in vino quam in aqua Spirituum copiam reperiri : ideoque fore ut si experimentum fieri posset in iis liquoribus, (fieri autem posse satis aperte negabant) Vinum plus spatii quam aqua relicturum esse in apice Tubi, positis Tubis eiusdem altitudinis.
\note{« cum » est écrit au-dessus de « de » biffé. Dans la marge de gauche, à hauteur de « autem pedes $31\frac{2}{3}$ », deux ou trois points à l'encre, sans lettre.}

Hoc concesso, ostensum est eis malum jacens cum Tubis alligatis : quibus altero aqua, altero vino repletis et oribus occlusis, erectum est Malum ; et Scutulæ oribus ipsis Tuborum admotæ, quarum altera vino, altera aqua plena erat, in quas immersa \struck{erant} erant Tuborum ora, Tubis semper plenis remanentibus, donec ora eorum recluderentur : quibus apertis, statim ambo illi liquores, in tubis contenti sic depressi sunt, ut postquam quiescerent, staret altitudo aquæ in suo tubo supra superficiem alterius aquæ in Scutula subiecta contentæ, pedibus $31\frac{1}{9}$ circiter, vinum autem paulo altius puta $31\frac{2}{3}$ proxime remanentibus Tuborum reliquis partibus \struck{\ill{}} superioribus, veluti Vacuis, omnino sicuti in hydrargyro deprehensum erat.
\note{« erant » est d'abord écrit puis biffé, et récrit aussitôt. Le mot biffé devant « superioribus » commence par une majuscule en A ; les lettres suivantes sont noyées dans la rature.}

Rursus autem mutati sunt liquores in Tubis, ut qui prius aqua, is postea vino repleretur, et vicissim, nec ideo quidquam in experimento mutatum est, circa vtramq; altitudinem. Inclinatum est malum sensim et sensim vterque Tubus replebatur ascendente liquore, ut semper obtineret prædictas altitudines secundum perpendicularem mensuratas. Facta est quoq; obseruatio in Tubo 15 pedum partim aqua, partim hydrargyro repleto : ac tum quoties illi duo liquores simul plus ponderabant quam altitudo pedum $2\frac{7}{24}$ solius hydrargyri ; aut quam altitudo pedum $31\frac{1}{9}$ solius aquæ quod eodem redit, toties posito in Scutella subiecta mero hydrargyro, ambo liquores in Tubo sic descenderunt, ut hydrargyrum inferiorem quidem Tubi
\note{« 15 » : le 5 est tracé comme un s long, de sorte que le nombre ressemble à « is » ; le même nombre revient à la vue 102 (« tubus 15 pedum »). « ac tum » est écrit d'un seul tenant. La phrase se poursuit à la vue 102.}

\page{102}
\note{Verso du feuillet numéroté 49, sans numéro. La même main continue la phrase de la vue 101 (« ut hydrargyrum inferiorem quidem Tubi / partem \ldots »).}

partem \struck{\ill{}} supra Scutellam occuparet. secundum altitudinem quæ minor esset pedibus $2\frac{7}{24}$. supra hydrargyrum extaret aqua, quæ altitudinem hydrargyro deficientem sic compensaret, ut ipsius aquæ altitudo contineret $13\frac{4}{7}$ partem altitudini hydrargyri deficientem : tum supra aquam reliquum tubi remaneret veluti Vacuum.
\note{le mot biffé après « partem » commence par une S majuscule ; il est récrit et raturé à plusieurs reprises, et n'a pas été lu. Le trait qui passe sur « 13 » prolonge la virgule qui suit « contineret » : le nombre n'est pas lu comme biffé, et il revient deux fois plus bas ($13\frac{4}{7}$). « altitudini » est écrit en fin de ligne, sans -s.}

Et quamquam ad hanc obseruationem adhibitus fuerit tunc tubus 15 pedum ; tamen illa postea in Tubis 4 aut 5 pedum facta est æque eleganter. Namq; ad hoc sufficit ut dum initio hydrargyrum in Tubum infunditur, illud non impleat quidem totum Tubum at idem obtineat altitudinem majorem pedibus $2\frac{7}{24}$ reliqua vero pars Tubi aqua repleatur sic enim res semper bene habebit : modo in Scutella subjecta merum sit hydrargyrum. Nam si huic imponatur aqua ; aut si Tubus aqua immergatur, siue totus, siue ex parte tantum ; tunc in omnibus obseruationibus siue merum hydrargyrum tubo includatur ; siue cum eo immittatur aqua, ambo illi liquores in Tubo altius ascendent, quam supra dictum est, et auctio illa continebit vnam partem, qualium altitudo aquæ hydrargiro Scutellæ superposita continebat $13\frac{4}{7}$ quandoquidem grauitas aquæ ad grauitatem hydrargyri sub aquali mole se habet, vt 1. ad $13.\frac{4}{7}$.
\note{« 15 » s'écrit ici encore avec un 5 en forme de s long. « aquali » est écrit ainsi, pour « æquali ». Dans le dernier « $13.\frac{4}{7}$ », le 4 est repassé.}

His et multo pluribus obseruationibus diuersimode multa cum diligentia exhibitis, et quater publice repetitis occlusum est os misellis illis Sciolis. Neque tamen ideo certi quidquam statui potuit circa spatium illud quod apparet veluti vacuum, in quo fit motus successiuus corporum et per quod transit lumen vna cum coloribus. Expectamus ergo dum Tubi fabricentur in \struck{quibus} \add{quos} quædam animalia intromittere possimus ut videamus an ibi viuere, ac præcipue an volare possint. Si enim volent procul dubio corpus est cui alæ innitentur, alias merum videbitur Vacu\ill{} ut ab eo permulti tantopere abhorrent ut quiduis potius quam tale spatium sine corpore admittant.
\note{« quos » est écrit au-dessus de « quibus » biffé. La ligne finit, au bord du texte, sur « Vacu » ; la suivante s'ouvre sur deux lettres (« vt » ou « ue » ?) avant « ut ab eo », qui n'ont pas été lues sûrement ; le sens attend « Vacuum ».}

Nec desunt qui maluerint ad materiam quandam subtilissimam recurrere in qua et aer, et cætera corpora tum dura quam fluida innatent, dum per eorum omnium poros illa facillime et nulla prope resistentia penetrat, prorsusq; omnia replet spatia. Quod quidem quia gratis affirmant, semper licebit gratis negare ; dum illi huiusmodi materiæ existentiam firmis rationibus stabilierint. Præsertim vero quia multis non difficilius videtur meri Vacui quam talis materiæ existentia cuius necessitas aut vsus in rerum natura nullus apparet.

Quidquid sit in eo jam fere omnes conueniunt, quod nisi merum illud sit Vacuum. at certe id eorum corporum vacuum est de quibus in Philosophia Peripatetica sermo habetur Puta, Terræ, Aquæ, aeris, Ignis et cœli : quodq; in liquore quod in obseruationibus adhibetur solius grauitatis non autem alterius qualitatis habenda sit ratio.
\note{ce paragraphe s'ouvre sur un Q initial plus grand, en retrait. Le « d » de « quod » (« in liquore quod ») est chargé d'une tache d'encre.}

Cæterum non omittendum est factam fuisse ab eodem D.\textsuperscript{no} de Paschal hydrargyri obseruationem in Tubis recuruis quos communiter Syphones vocamus : quibus hydrargyro repletis, tum ita inuersis, ut ambo orificia in hydrargyrum duabus Scutellis contentum immergerentur, Ventre Syphonis in aere extante versus superiorem partem : Siquidem Venter ille supra hydrargyrum Scutellæ \struck{\ill{} \uncertain{ex hydrargyro} \ill{}} \add{vtriusque minus ascenderet quam pedibus} $2\frac{7}{24}$ perpendiculariter sumptis, essentque
\note{la dernière ligne a été corrigée : trois ou quatre mots biffés d'un trait continu, dont le premier commence par un c et le dernier paraît finir en -ens ou -ere, sont remplacés par « vtriusque minus ascenderet quam pedibus », écrit au-dessus. « Syphonis » : la fin du mot, en bout de ligne, est serrée contre le bord du texte. La phrase se poursuit à la vue 103.}

\page{103}
\note{Recto. En haut à droite, à l'encre, « 50 » (série continue ; voir l'en-tête). La même main continue la phrase de la vue 102 (« essentque / Scutellæ cum suo hydrargyro \ldots »). Le cachet rond de la bibliothèque est apposé sur les dernières lignes, sans en cacher la lecture.}

Scutella cum suo hydrargyro secundum eandem superficiem horizontalem dispositæ, habita ratione superficiei superioris vtriusque hydrargyri in Scutellis contenti, ad horizontem, in vnam superficiem librarentur. At si venter Syphonis supra vtriusque Scutellæ hydrargyrum altior esset quam pedibus $2\frac{7}{24}$ perpend. tunc, (quod minime expertis miraculum videbatur ex vtroque crure Syphonis, hydrargyrum in Scutellam suam eo vsque recidebat, quousq; remaneret altitudo toties jam dicta pedum $2\frac{7}{24}$ supra hydrargyrum vtriusque Scutellæ siue ipsum in vtraque Scutella, in eodem plane horizontali staret siue non, Venter autem Syphonis apparebat veluti vacuus inter duas partes hydrargyri hinc inde in cruribus Syphonis contentas : quo in statu cessabat omnis motus. Tandem si venter Syphonis supra vnius Scutellæ hydrargyrum minus altus esset alterius vero vel æque altus vel altior mensura prædicta, tunc pro varia altitudinis differentia variæ phases conspiciebantur quas recensere nimis longum esset.
\note{« perpend. » est abrégé ainsi. La parenthèse ouverte avant « quod minime expertis » n'est pas refermée sur la page.}

Nec dubium fuit quin idem omnino in aqua, vel in alio quouis liquore fieret si Syphones debitæ altitudinis adhiberentur. Ne quis ergo putet in posterum aquam ex vno latere montis quantumuis alti posse, superato vertice montis, firmiss.\textsuperscript{mi} Syphonis beneficio deduci in Vallem opposito montis lateri subjacentem : hoc enim supra altitudinem \struck{\ill{}} pedum 31 fieri non potest. Neque etiam supra hanc altitudinem ascendet aqua in tubis illis qui communiter puteis applicantur ad suggendam et attollendam aquam. si machina adhibita suggat tantum et non premat : nam premendo poterit sane aqua altius pelli, si machina debite constructa sit ac disposita.
\note{« firmiss.\textsuperscript{mi} » est écrit dans la marge de gauche, en tête de la ligne, et se lie à « Syphonis » : c'est un ajout de la même main, qui se lit dans la phrase. Le mot biffé après « altitudinem » commence par a- et paraît être une forme de « ascendere » ; il n'a pas été lu sûrement. Le nombre est ici « 31 » seul, sans fraction.}

His Animaduersis inter multas quæstiones duæ nobilissimæ inter eruditos agitari cœperunt. Prima. Qui fieret ut liquor qui in obseruatione adhibetur in certa quadam altitudine supra subjectum sibi in Scutella liquorem veluti inter vires quasdam contrarias libratus sustineretur habita ratione ponderis ipsius liquoris ad altitudinem ipsius in Tubo determinandam ?. Secunda Qui fieret ut quiescente liquore in debita sibi altitudine si vel gutta ab eo \struck{descendat} \add{decedat}, admissa scilicet per os Tubi quadam cuiusuis alius liquoris gutta, tum reliquus liquor in Tubo contentus, ad summum Tubi subito, et, si hydrargyrum fuerit, magno cum strepitu ascendat, et tanta vi in ipsum fundum superius impingat ut multos Tubos præcipue ex longioribus fregerit, atque in aerem longe supra Tubum exilierit : si que Tubus illæsus supersit, rursus descendat liquor, et per aliquam motuum contrariorum reciprocationem secundum ascensum et descensum, quæ reciprocatio omnibus ferme experimentis familiaris est quiescat : tum ; dimissa alia gutta ascendat rursus, minus tamen violenter : atque ita alternando per aliquot vices donec totus liquor grauior ex Tubo dimissus fuerit, et leuior in locum illius ascenderit !.
\note{« His » ouvre un nouvel alinéa, en retrait, avec une grande majuscule. « decedat » est écrit au-dessus de « descendat » biffé. La première question se ferme sur « ?. », la seconde sur un signe qui ressemble à « !. » ; les deux sont transcrits tels qu'ils paraissent.}

Sed quia circa has quæstiones eruditi diuersa sentiunt et discussio opinionum longa esset ideo expectabo super hac re Epistolas tuas, ac tunc si ita videbitur vltra progrediar et si quid interea noui acciderit quod tibi placiturum putem illud narrationi adjungam. Vale Vir mihi Colendiss.\textsuperscript{me} et me ut hucusque fecisti in posterum amare perge.
\note{« hac re » est écrit d'un seul tenant. Le cachet de la bibliothèque couvre en partie « esset ideo » et la fin de « adjungam », qui se lisent au travers.}

Parisijs 12 Calend. Octob. 1647
\note{date seule, en retrait, au bas de la page, le millésime souligné d'un long trait. « Calend. » et « Octob. » sont abrégés, avec un trait suscrit. Le dernier chiffre du millésime a la forme à crochet que cette main donne au 7 (ainsi le numérateur de $2\frac{7}{24}$ à la première ligne de la vue 101, dont la valeur est fixée par les nombres du texte) ; il diffère nettement du 4 fermé qui le précède. Aucune souscription ni signature : la lettre, qui s'adresse à un correspondant à la deuxième personne (« Epistolas tuas », « tibi », « Vale Vir mihi Colendiss.\textsuperscript{me} »), ne nomme ni son auteur ni son destinataire sur ces pages.}

\end{document}
